\section{Motivation and significance}
\label{sec:intro}

Mass phishing is, for practical purposes, a solved filtering problem: modern email security
intercepts the overwhelming majority of bulk campaigns. The residual, disproportionately
costly threat is \emph{targeted}---spear phishing, business email compromise (BEC), and
lateral phishing sent from an already-compromised internal account
\cite{ho2019detecting,cidon2019high,apwg2024trends}. These messages are hard precisely
because they are \emph{contextually legitimate}: the sender is a plausible or genuine
contact, the timing is unremarkable, and the request often continues a real thread. They
carry no malware and little lexical anomaly, so content-only classifiers degrade toward
chance on exactly the cases that inflict the largest losses.

A natural response is to stop looking only at the message and start looking at the
\emph{relations} around it. If one models an organization's communication as a graph---who
talks to whom, when, and through which channel---then a targeted attack becomes a structural
or temporal irregularity: an edge that should not exist, a sender whose neighbourhood
behaviour shifts, a burst that violates the local rhythm. This intuition has driven a decade
of work, from account-compromise heuristics \cite{egele2013compa,stringhini2015thataint}
and lateral-phishing feature models \cite{ho2019detecting} to graph neural networks (GNNs)
\cite{kipf2017semisupervised,hamilton2017inductive} adapted from fraud and anomaly detection
\cite{dou2020enhancing,ma2021comprehensive}.

\paragraph{The problem: evaluation, not models.} What the field lacks is not another model
but a way to \emph{believe} the numbers. Real lateral-phishing corpora are proprietary,
small, and saturated with confounders. Attacker events systematically differ from benign
events in send-time rhythm (off-hours activity), sender degree (unusual recipient sets),
volume, and identity (a few repeated malicious senders). A detector can therefore achieve a
high score by keying on these artefacts rather than on the mechanism of the attack---and
such a detector fails on a new tenant or after the attacker adapts. This is the security-ML
measurement trap documented broadly by \citet{arp2022dos} and quantified for temporal drift
by \citet{pendlebury2019tesseract}. In the graph setting the risk is acute: node degree and
timing are exactly the features a GNN will happily overfit. No open, reusable testbed today
offers, simultaneously, (i) ground-truth labels, (ii) a controllable \emph{adaptive}
attacker, and (iii) explicit leakage controls. As a result, published results are point AUCs
on incomparable private data, and it is impossible to tell mechanism from confounder.

\paragraph{Contributions.} We address this with \tool, an installable, extensible toolkit
for adversarial, leakage-controlled evaluation of graph-based phishing and lateral-movement
detectors. Concretely:

\begin{itemize}[leftmargin=1.4em,itemsep=2pt]
\item \textbf{A controllable testbed} coupling three graph sources (a procedural
  organization model, SNAP email networks \cite{leskovec2014snap}, and the Enron corpus
  \cite{klimt2004enron}) with a programmable \emph{cascade} attacker whose reach/evasion
  trade-off is a small set of parameters, producing scenarios with per-event ground truth
  (Section~\ref{sec:method}).
\item \textbf{A leak-aware evaluation protocol}---benign traffic matched by design to the
  attack's endpoint and degree distribution, victim-disjoint splits, a time-shuffle
  causality control, and operational Recall@low-FPR metrics---that separates the attack
  mechanism from the four confounders above (Section~\ref{sec:method}).
\item \textbf{A detector panel behind one interface}, spanning tabular baselines,
  account-compromise heuristics, an anomaly forest, and static/temporal GNNs, so that
  heterogeneous methods are compared on identical, audited footing.
\item \textbf{An automatic red/blue arms-race} that co-searches attack and defense spaces to
  produce robustness frontiers rather than single-point scores.
\item \textbf{An open, GDPR-free artefact}: the synthetic population contains no personal
  data, and all experiments reproduce offline from bundled result files.
\end{itemize}

\paragraph{Analyses the tool supports.} \tool operationalizes analyses that closed corpora
cannot, which the authors' companion studies on graph-based spear-phishing detection (in
preparation) rely on: \textbf{(A1)} which graph structures carry signal for which attack type;
\textbf{(A2)} whether a detector's advantage survives a time-shuffle---mechanism or leakage;
\textbf{(A3)} how detection degrades as an attacker trades reach for evasion; and
\textbf{(A4)} whether a detector's robustness transfers across conditions. This paper describes
the \emph{software} that makes such analyses reproducible; the corresponding scientific results
and their interpretation are reported separately. Section~\ref{sec:related} positions the tool
against detection methods and evaluation critiques; Section~\ref{sec:method} describes the
software; Section~\ref{sec:results} gives illustrative examples;
Section~\ref{sec:discussion} details impact and limitations;
Section~\ref{sec:conclusion} concludes.